\chapter{The Periodic Table Is Not an Element Directory}

\begin{kaichang}
Open a drawer at home and you will probably find a battery, aluminum foil, a stainless-steel spoon, a packet of salt, and a nail. Look outside: leaves are green, the sky is blue, and your exhaled breath contains carbon dioxide.

Now play a game. Which raw materials do these objects share? Foil and drink cans both contain aluminum; that is easy. But are the iron in a steel spoon, the iron that makes blood red, and the iron in a rusty nail the same iron? Is chlorine in table salt the same thing as chlorine in swimming-pool disinfectant? What about magnesium in green leaves and magnesium supplements?

Write down your guesses. At the end, we will check them using a table you have probably seen, but perhaps never really understood: the periodic table.
\end{kaichang}

\section{A Directory Hanging in the Wrong Place}

Many people's first encounter with the periodic table is on a chemistry classroom wall: over a hundred colored boxes, each with a symbol and number. It looks like a directory: hydrogen at number 1, helium 2, oxygen 8, gold 79. If that were all, learning chemistry would mean simply memorizing it.

But examine its shape. This is no ordinary rectangle. A corner is missing on the right, with two detached rows below, like an oddly shaped city. Why does the first row have only 2 boxes, the second and third 8 each, then 18 and 32? Why are the bottom lanthanide and actinide rows exiled outside?

A directory would need none of these shapes; it could be rectangular. The strange periodic table is actually nature's own map, telling you that \textbf{elemental personalities recur in cycles}.

What does recurring personality mean? Consider lithium (\chem{Li}), sodium (\chem{Na}), and potassium (\chem{K}). All are silvery metals soft enough to cut with a knife, stored in kerosene, and violently reactive with water, with increasing temper: lithium hisses and circles, sodium melts into a darting bead, and potassium bursts into violet flame. Three elements with matching personalities, placed in one column.

Now the far-right column: helium, neon, argon, krypton, xenon, radon. This family's personality is the opposite: reluctant to do anything. Helium in balloons, neon in signs, and argon protecting welds avoid bonding with other elements and keep to themselves.

One column, one personality: the table's greatest secret. But why? Elements cannot convene to agree. The answer lies inside atoms, in Chapter 9's electrons, rather than in the table itself.

\section{Electron Floors Decide Everything}

Recall Chapter 9: electrons do not fly about randomly outside the nucleus. They occupy energy floors with strict capacities: 2 on the first, 8 on the second, 18 on the third, and so on. They fill upward from lower energies, like hotel guests starting on the ground floor.

Chemical personality---readiness to react, preferred partners, and whether electrons are surrendered or seized---depends almost entirely on \textbf{how many electrons occupy the outermost floor}. Inner electrons are thoroughly screened by the outer ones and seldom appear.

Read the table horizontally, and something remarkable emerges:

\begin{itemize}[itemsep=2pt]
\item Hydrogen and helium fill only the first shell's 2 places, so row one has 2 boxes.
\item Lithium through neon fill the second shell's 8 places, giving row two exactly 8 boxes.
\item Sodium through argon fill 8 outer places in the third shell, so row three also has 8.
\item Farther down, greater shell capacities widen the table to 18 and 32 boxes.
\end{itemize}

Vertically, elements in one column have equal outer-electron counts. Lithium, sodium, and potassium each have one lonely outer electron and readily surrender it, sharing fiery tempers. Helium, neon, and argon have full outer shells and want nothing. Widths, row counts, and family resemblances all arise from electron lodging rules.

\begin{wuqu}
``Elements react to complete an octet of 8 electrons.'' This common misconception assigns atoms desires they do not have. The real statement is that a change may occur when giving, taking, or sharing electrons lowers the entire system's energy. Eight outer electrons often correspond to a lower-energy arrangement, as in noble gases, making the octet a useful rule of thumb, not a natural law. Many stable molecules and ions violate it, as later chapters repeatedly show (Chapters 17 and 19). Treat the mnemonic as shorthand for energy rules, not an atom's purpose in life, and you will not be misled.
\end{wuqu}

\section{Mendeleev's Great Gamble}

\begin{zhengju}
In 1869, the Russian chemist Mendeleev arranged dozens of cards on his desk, each describing a known element and its properties. Only 63 elements were known, with no concepts of electrons or nuclei; the electron would not be discovered until 1897. His sole clue was that ordering by increasing atomic weight made properties \textbf{repeat periodically}.

He did something that seemed almost reckless. First, he left gaps, declaring undiscovered elements and predicting their densities, melting points, and chemistry from neighboring personalities. He named one gap eka-aluminum. Second, where weight order conflicted with personality, as for tellurium and iodine, he insisted on personality, preferring to doubt measured weights.

Sixteen years later, in 1875, the French chemist Boisbaudran discovered gallium (\chem{Ga}), whose position, density, and properties strikingly matched eka-aluminum. Germanium (\chem{Ge}), discovered in 1886, fulfilled eka-silicon's prediction. A table predicting unseen things is science at its most confident. Tellurium and iodine were eventually resolved too: Mendeleev's personality order was correct because nuclear proton count, atomic number, sets the true order, not atomic weight. This was Chapter 7's elemental identity card.

Remember the method rather than simply Mendeleev's cleverness. He trusted that repeated patterns had a shared cause and made that trust falsifiable through empty boxes. Evidence arrived and the pattern held, though he could not yet explain why. Quantum mechanics supplied the reason, electron configurations, half a century later.
\end{zhengju}

\section{A Map of the Table's Neighborhoods}

Tour the modern periodic table as a city. A few neighborhoods keep you oriented:

\begin{itemize}[itemsep=2pt]
\item \textbf{The leftmost two columns}, Groups 1 and 2, alkali and alkaline-earth metals: only 1--2 outer electrons, eagerly surrendered. These active neighborhoods readily react with water and acids.
\item \textbf{The broad low-rise central district}, transition metals: iron, copper, zinc, titanium, chromium, and manganese live here. Steady yet versatile, they show several valences and often colorful ions (Chapter 15 tours the district).
\item \textbf{The second column from the right}, Group 17, halogens: fluorine, chlorine, bromine, iodine have 7 outer electrons, one short of full. They seize electrons readily and become stable negative ions, as chlorine does to form chloride in salt.
\item \textbf{The far-right column}, Group 18, noble gases: full outer shells and little interest in anything.
\item \textbf{The detached lanthanides and actinides below}: their differences occur in deeper inner electrons while outer shells remain nearly identical, so 14 elements have extremely similar personalities. Inserting them into the main table would make it too wide to print, hence the downstairs annex.
\end{itemize}

An invisible diagonal runs from boron toward astatine. Metals to its left conduct, deform, and readily lose electrons; nonmetals lie rightward. Borderline metalloids such as silicon and germanium share features of both. Silicon's half-metal, half-nonmetal personality makes your phone chip work.

\begin{qiaoliang}
Why does reactivity vary regularly down a group? Lithium, sodium, and potassium all lose one outer electron, so why is potassium fiercest? Distance and screening matter. Potassium's outer electron occupies shell four, farther from the nucleus behind three inner-electron blankets that screen attraction (Chapter 11 explains shielding). Loosely held electrons leave readily; easier loss makes the water reaction fiercer. In Chapter 27's energy language, net released energy and the energy needed to start both vary with shell height. Every trend arrow reflects a struggle between nuclear attraction and electron repulsion and screening. Understand the struggle rather than memorize arrows.
\end{qiaoliang}

\section{Back to the Opening Drawer}

Now we can keep our promise. Iron in a steel spoon, blood, and rust is the same element, \chem{Fe}, from the central neighborhood. In the spoon it is metallic atoms; in blood's hemoglobin, a coordination center binding oxygen (Chapter 15); in rust, an oxidized form that has surrendered electrons to oxygen. Different electronic states and neighbors give one element entirely different roles.

Chlorine in salt and disinfectant is likewise one element, but chloride \chem{Cl^-} has acquired an electron and settled down, while reactive chlorine molecules \chem{Cl_2} have not. Unchanged proton-count identity coexists with utterly different personalities and hazards, which is why eating salt and inhaling chlorine are completely different.

Magnesium in green leaves and supplements is the same Group 2 element. In a leaf, it sits at the exact center of chlorophyll, a crucial screw in photosynthesis machinery (Chapter 54).

Elements never appear in everyday life as pure personalities; they always occupy particular combinations and electronic states. The table supplies the underlying cards for understanding substances, not a simple list of poisonous versus harmless things.

\section{Our City Tour Is Not Finished}

We have seen only the map's outline. Chapter 11 unpacks atomic radius, ionization energy, and electronegativity trends from the electron-floor struggle. Chapters 12--15 visit fiery Group 1, electron-hungry Group 17, life's structural elements oxygen, sulfur, nitrogen, and phosphorus, and colorful transition metals. The map's boundaries also keep expanding: newly synthesized elements are still appended, with floors high enough to challenge familiar rules. This is evidence that science remains alive.

\begin{tiaozhan}
Why 8 places on shell two and 18 on shell three? These are not arbitrary rules. Quantum mechanics supplies four quantum numbers: principal $n$ for floor, angular $l$ for room type s, p, d, f, magnetic for orientation, and spin for at most two opposite-spin electrons per room under Pauli exclusion. Shell $n$ holds $2n^2$: $n=1$ gives 2, $n=2$ gives 8, $n=3$ gives 18. Period lengths directly reflect these numbers. Filling order has one twist: fourth-shell s rooms are cheaper, lower in energy, than third-shell d rooms, so electrons occupy 4s before returning to 3d. This creates the fourth row's 10 extra transition-metal boxes and underlies chromium's and copper's irregular configurations. Chapter 9's challenge and Chapter 11 develop the details.
\end{tiaozhan}

\section*{Try It Yourself}

\begin{sikaoti}
\item Without consulting the table, use outer-electron counts to predict which of magnesium (\chem{Mg}, 2 outer electrons) and oxygen (\chem{O}, 6) will give electrons and which take them. What charges will they have in the product?
\item Helium's first shell is full with 2 electrons; hydrogen has only 1. Explain from shell rules why hydrogen burns while helium safely fills balloons.
\item Why did Mendeleev put tellurium before iodine despite its greater atomic weight? What later discovery vindicated him?
\item Draw sodium's 11 and chlorine's 17 electrons as floors and rooms. Mark which outer floor is almost full and which has just opened. Note that electrons do not really live in little boxes; this representation helps us reason.
\item Fluorine tops Group 17, with chlorine below. Reverse Group 1's logic that higher shells lose electrons more easily: how should electron-taking ability change upward within a group? Consult the periodic table to check.
\end{sikaoti}
